\documentclass{article}
\usepackage[margin=1in]{geometry}
\pagestyle{empty}
\begin{document}
\section*{Sales datamart: transformed tables}
\noindent Entity--relationship diagram for the tables emitted by the Beam
normalization graph. The same schema is used by the local report checks and
the configured BigQuery sink. PK means primary key; FK means foreign key.

\begin{center}
\setlength{\unitlength}{1pt}
\begin{picture}(450,330)
\put(0,220){\framebox(205,105){}}
\put(0,299){\line(1,0){205}}
\put(0,299){\makebox(205,26){\textbf{customers\_Dataflow}}}
\put(9,229){\parbox[b][62pt][t]{190pt}{\small
  \begin{tabular}{@{}lll@{}}
    \texttt{customer\_id} & STRING & PK \\
    \texttt{name} & STRING & \\
    \texttt{region} & STRING &
  \end{tabular}}}

\put(245,220){\framebox(205,105){}}
\put(245,299){\line(1,0){205}}
\put(245,299){\makebox(205,26){\textbf{products\_Dataflow}}}
\put(254,229){\parbox[b][62pt][t]{190pt}{\small
  \begin{tabular}{@{}lll@{}}
    \texttt{product\_id} & STRING & PK \\
    \texttt{name} & STRING & \\
    \texttt{category} & STRING & \\
    \texttt{price\_cents} & INT64 &
  \end{tabular}}}

\put(112,0){\framebox(225,140){}}
\put(112,114){\line(1,0){225}}
\put(112,114){\makebox(225,26){\textbf{orders\_Dataflow}}}
\put(121,9){\parbox[b][97pt][t]{208pt}{\small
  \begin{tabular}{@{}lll@{}}
    \texttt{order\_id} & STRING & PK \\
    \texttt{customer\_id} & STRING & FK \\
    \texttt{product\_id} & STRING & FK \\
    \texttt{order\_date} & DATE & \\
    \texttt{quantity} & INT64 &
  \end{tabular}}}

\put(102,220){\line(0,-1){44}}
\put(102,176){\line(1,0){65}}
\put(167,176){\vector(0,-1){36}}
\put(111,201){\makebox(32,14)[l]{1}}
\put(173,147){\makebox(36,14)[l]{0..*}}
\put(112,180){\makebox(75,14)[l]{\small places}}

\put(347,220){\line(0,-1){44}}
\put(282,176){\line(1,0){65}}
\put(282,176){\vector(0,-1){36}}
\put(320,201){\makebox(18,14)[r]{1}}
\put(242,147){\makebox(34,14)[r]{0..*}}
\put(263,180){\makebox(75,14)[r]{\small appears in}}
\end{picture}
\end{center}

Each customer can place many orders; each product can appear in many orders.
Each order references exactly one customer and one product. The arrows point
from each parent to its child table. A parent can have zero or many orders.

\paragraph{Integrity boundary.} These are logical relationships checked by
\texttt{reports.py} and \texttt{sql/bigquery.sql}, not automatically enforced
BigQuery foreign-key constraints. The fixture intentionally contains one
orphan customer reference; the check reports it, and inner-join reports
exclude that order. Duplicate or conflicting source keys are resolved or
rejected before this reporting stage.

\paragraph{Money.} \texttt{price\_cents} is the normalized integer-cent value,
not the original raw price string. Revenue multiplies it by the order
quantity with explicit numeric-range checks. No cloud execution is implied
by this schema diagram.
\end{document}
